%!TEX program = xelatex 
\documentclass[a4paper,12pt]{article}
\usepackage{fontspec}\defaultfontfeatures{Ligatures=TeX}
\usepackage[a4paper,vmargin={3cm,3.5cm},hmargin={3cm,3cm}]{geometry}
%-----------------------------------------------------------------------------%
\usepackage{settings}
\usepackage{tikz}
%-----------------------------------------------------------------------------%
%%% Title %%%
\title{G6411 Recitation Notes: Hypothesis Testing: Basics}
\author{Jesse Chieh Chen\\\href{mailto:cc5229@columbia.edu}{cc5229@columbia.edu}}
\date{\today}
%-----------------------------------------------------------------------------%

\begin{document}

\maketitle

\section{Hypotheses and Tests: Terminology}

\subsection{What is a Hypothesis?}

A statistical model specifies a collection of possible distributions for the data.
A \emph{hypothesis} restricts this collection.
It may specify a whole distribution, a parameter value, or a property such as a conditional mean.
Write $\theta\in\Theta$ for the unknown features indexing the model and consider
\begin{align}
	H_0:\theta\in\Theta_0
	\qquad\text{against}\qquad
	H_1:\theta\in\Theta_1,
\end{align}
where $\Theta_0$ and $\Theta_1$ are disjoint.
Usually $\Theta_1=\Theta\setminus\Theta_0$.
The \emph{null hypothesis} $H_0$ is the restriction we test;
the \emph{alternative hypothesis} $H_1$ describes the departures.

\begin{definition}[Simple and Composite Hypotheses]
	A hypothesis is \emph{simple} if it specifies a single distribution for the data.
	It is \emph{composite} if it allows more than one distribution.
\end{definition}

\begin{example}
	Suppose $Y_1,...,Y_n$ are independent $\normaldistribution(\mu,1)$ variables.
	The hypothesis $H_0:\mu=0$ is simple,
	whereas the hypothesis $H_0:\mu\le0$ is composite.
\end{example}

\subsection{Decision Rules and Test Statistics}

\begin{definition}[Test and Rejection Region]
	Let $D_n$ denote the full sample.
	A test is a decision rule $\varphi_n(D_n)\in\{0,1\}$,
	where one means reject $H_0$ and zero means do not reject $H_0$.
	Its \emph{rejection region} is the set of samples
	\begin{align}
		\mathcal R_n=\{d:\varphi_n(d)=1\}.
	\end{align}
\end{definition}

\begin{definition}[Test Statistic]
	A test statistic is a function of the data $T_n(D_n)$ used to construct a test.
\end{definition}

\begin{remark}[Why use Test Statistics]
	We often construct a test using a scalar test statistic.
	For example, we might reject when $T_n>c$ or when $|T_n|>c$.
	This reduces a potentially high-dimensional sample to a quantity whose distribution we can study.
	A useful statistic should have a tractable distribution under the null
	and tend to enter the rejection region under alternatives.
\end{remark}

\begin{remark}[Cannot Reject]
	Not rejecting $H_0$ does not establish that it is true.
	The test may simply be unable to distinguish the null from the relevant alternatives with the available data.
	Nevertheless, in common parlance we often use ``accept'' to mean ``cannot reject.''
\end{remark}

\section{Errors, Size, and Power}

\begin{table}[t]
	\centering
	\begin{tabular}{l|cc}
		\toprule
		& $H_0$ is true & $H_1$ is true \\
		\midrule
		Reject $H_0$ & Type I, False Positive &  \\
		Accept $H_0$ &  & Type II, False Negative \\
		\bottomrule
	\end{tabular}
	\caption{Type I and Type II Errors.}
	\label{tab-errors}
\end{table}

As shown in \autoref{tab-errors}, a Type I error occurs when we reject a true null hypothesis,
i.e., a false positive.
A Type II error occurs when we fail to reject the null even though the alternative is true,
i.e., a false negative.

\begin{definition}[Power, Size, and Level]
	The \emph{power function} is the rejection probability at each parameter value:
	\begin{align}
		\pi_n(\theta)=\Pr_\theta\{\varphi_n(D_n)=1\}.
	\end{align}
	Under the null, this is the Type I error probability.
	Under an alternative, it is power, and $1-\pi_n(\theta)$ is the Type II error probability.
	The \emph{size} of the test is
	\begin{align}
		\sup_{\theta\in\Theta_0}\pi_n(\theta).
	\end{align}
	A test has \emph{level} $\alpha$ if its size is at most $\alpha$.
\end{definition}

\begin{remark}
	The supremum matters for a composite null: controlling the rejection probability at just one null value is not enough.
	The significance level $\alpha$ is chosen before looking at the test result.
	A common choice is $\alpha=0.05$.
\end{remark}

\begin{remark}[The Neyman-Pearson Approach]
	There is a trade-off between the two errors.
	Never rejecting eliminates Type I errors but has zero power;
	always rejecting has power one but size one.
	The Neyman-Pearson approach fixes a bound $\alpha$ on size
	and seeks high power subject to that bound.
	Different tests may have higher power against different alternatives,
	so a test that is best against every alternative need not exist.
\end{remark}

\begin{definition}[Consistent Test]
	A sequence of tests is \emph{consistent} against $\Theta_1$ if,
	for every fixed $\theta\in\Theta_1$,
	\begin{align}
		\pi_n(\theta)\to1\qquad\text{as }n\to\infty.
	\end{align}
\end{definition}

\begin{remark}
	Consistency concerns power, not size; a useful test needs both size control and good power.
	The word \emph{fixed} is important: consistency does not promise power approaching one
	against alternatives that move toward the null as the sample grows.
\end{remark}

\section{Critical Values and \texorpdfstring{$p$}{p}-Values}

\subsection{Critical Values}

\begin{definition}[Critical Value]
	A \emph{critical value} is a cutoff for a scalar test statistic $T_n$
	that is used to define the rejection region.
\end{definition}

\begin{remark}
	Typically, it is selected using a null distribution.
	If large values of $T_n$ lead to rejection and its continuous null \gls{cdf} is $G_0$,
	the rejection region can be defined as $\{T_n>c_{1-\alpha}\}$
	where $c_{1-\alpha}=G_0\inv(1-\alpha)$ is the $(1-\alpha)$-quantile of the null distribution.
	Then the test has level $\alpha$:
	\begin{align}
		\Pr_{H_0}(T_n>c_{1-\alpha})=\alpha.
	\end{align}
\end{remark}

\begin{example}
	For a two-sided statistic with standard normal null distribution,
	write $\Phi$ for the standard normal \gls{cdf} and $z_q=\Phi\inv(q)$.
	Then reject when
	\begin{align}\label{eq-normal-rejection}
		|T_n|>z_{1-\alpha/2}.
	\end{align}
	The two tails each have probability $\alpha/2$.
	At the 5\% level, $z_{0.975}\approx1.96$.
	For an upper-sided alternative, reject when $T_n>z_{1-\alpha}$;
	for a lower-sided alternative, reject when $T_n<z_\alpha$.
	The direction of the test should be chosen before seeing the estimate.
\end{example}

\subsection{The Observed Tail Probability}

\begin{definition}[$p$-value]
	For a specified test statistic and ordering of what counts as more extreme,
	the $p$-value is the null probability of a statistic \emph{at least as extreme as the observed value}.
\end{definition}

\begin{remark}
	Instead of looking at whether $T_n$ falls into a pre-specified rejection region,
	we can simply compare the $p$-value implied by the observed statistic to the significance level $\alpha$.
	If $p<\alpha$, we reject; if $p\ge\alpha$, we do not reject.
\end{remark}

\begin{example}
	For a two-sided standard normal test statistic $T_n$ with realized value $t_{\mathrm{obs}}$,
	the $p$-value is defined as
	\begin{align}\label{eq-normal-pvalue}
		p
		\coloneqq \Pr\big(|T_n|\ge|t_{\mathrm{obs}}|\big)
		= 2\big(1-\Phi(|t_{\mathrm{obs}}|)\big).
	\end{align}
	For an upper-sided test, $p\coloneqq1-\Phi(t_{\mathrm{obs}})$;
	for a lower-sided test, $p\coloneqq\Phi(t_{\mathrm{obs}})$.
\end{example}

\begin{remark}[Validity of $p$-values]
	A valid $p$-value satisfies $\Pr_\theta(p\le a)\le a$ for every $\theta\in\Theta_0$ and $0<a<1$.
	Under any parameter satisfying the null, the probability of obtaining a $p$-value at most $a$ is at most $a$.
	Thus rejecting when $p\le\alpha$ controls the Type I error probability at level $\alpha$.
\end{remark}

\section{Normal-Based Testing}

\subsection{Pivots and Asymptotic Pivots}

\begin{definition}[Pivotal Quantity]
	A pivotal quantity is a function of the data and parameters whose distribution does not depend on unknown parameters.
\end{definition}

\begin{example}
	For independent normal observations with mean $\mu$ and known variance $\sigma^2$,
	$\sqrt n(\bar Y_n-\mu)/\sigma\sim\normaldistribution(0,1)$ is pivotal
	since standard normal distribution does not depend on $\mu$ and $\sigma$.
\end{example}

\begin{remark}
	Exact pivots typically rely on strong distributional assumptions.
	Our main approach instead uses an \emph{asymptotic pivot}.
	That is, we construct a statistic whose \emph{limiting} distribution does not depend on unknown parameters.
\end{remark}

\begin{example}[Asymptotic Normal Testing]
	Let an estimator $\widehat\theta\in\reals^k$ satisfy the following, with $V$ positive definite:
	\begin{align}
		\sqrt n(\widehat\theta-\theta)\dto\normaldistribution(0,V)
		\quad\text{which implies}\quad
		V^{-1/2}\sqrt n(\widehat\theta-\theta)\dto\normaldistribution(0,I_k).
	\end{align}
	Suppose that we have a consistent variance estimator $\widehat V\pto V$.
	Under $H_0:\theta=\theta_0$,
	construct the statistic $T_n$ as follows
	\begin{align}\label{eq-studentized-statistic}
		T_n
		\coloneqq \widehat V^{-1/2} \sqrt n(\widehat\theta-\theta_0)
		= \Big(\widehat V^{-1/2} V^{1/2}\Big) V^{-1/2} \sqrt n(\widehat\theta-\theta_0)
		\dto\normaldistribution(0,I_k)
	\end{align}
	where the asymptotic distribution of $T_n$ follows from Slutsky's theorem.
	Hence, $T_n$ is asymptotically pivotal.
\end{example}

\begin{remark}[Pointwise Asymptotic Validity]
	Choosing a rejection region with probability $\alpha$ under the limiting normal distribution
	gives rejection probability approaching $\alpha$ at each fixed null parameter where the stated assumptions hold,
	provided the region's boundary has probability zero under that limiting distribution.
	For a composite null, a uniform bound on size over all null parameters requires additional uniformity conditions.
\end{remark}

\subsection{Testing a Regression Coefficient}

Consider $Y_i=X_i\T\beta+e_i$ with $X_i\in\reals^k$ and $\E(X_ie_i)=0$.
Suppose the observations are \gls{iid} and the rank and moment conditions for
\gls{ols} asymptotic normality and consistent variance estimation hold.
Define $g_i\coloneqq X_i e_i$, $Q_{xx}\coloneqq\E X_iX_i\T$, and $\Omega\coloneqq\E g_i g_i\T$.
Then we have
\begin{align}
	\sqrt n(\widehat\beta-\beta)\dto\normaldistribution(0,V)
	\quad\text{where}\quad
	V \coloneqq Q_{xx}\inv\Omega Q_{xx}\inv.
\end{align}
Suppose that we want to test a slope restriction $H_0:\beta_j=b_{j0}$ against $H_1:\beta_j\ne b_{j0}$.
Previously, we have constructed a consistent estimator $\widehat V$ of $V$ using the sample analogs as follows
\begin{align}
	\widehat V
	&\coloneqq
	\left(\frac1n\sum_{i=1}^{n}X_iX_i\T\right)\inv
	\left(\frac1n\sum_{i=1}^{n} \widehat g_i\widehat g_i\T\right)
	\left(\frac1n\sum_{i=1}^{n}X_iX_i\T\right)\inv
	\pto
	V
\end{align}
where $\widehat g_i\coloneqq X_i\widehat e_i$ and $\widehat e_i=Y_i-X_i\T\widehat\beta$.
Under $H_0:\beta_j=b_{j0}$, selecting the $j$-th coordinate gives
\begin{align}
	\sqrt n(\widehat\beta_j-b_{j0})\dto\normaldistribution(0,V_{jj}).
\end{align}
Provided $V_{jj}>0$, Slutsky's theorem gives
\begin{align}
	T_n=\widehat V_{jj}^{-1/2}\sqrt{n}(\widehat\beta_j-b_{j0})
	=\frac{\widehat\beta_j-b_{j0}}{\operatorname{se}(\widehat\beta_j)}
	\dto\normaldistribution(0,1)
	\quad\text{where}\quad
	\operatorname{se}(\widehat\beta_j)=\sqrt{\frac{\widehat V_{jj}}{n}}.
\end{align}
The approximate two-sided $p$-value is $p=2(1-\Phi(|T_n|))$ as given in \eqref{eq-normal-pvalue}.

\begin{remark}[Size and Power]
	Let's check the size and power of this test.
	Under the null, $T_n\dto\normaldistribution(0,1)$,
	so the rejection probability approaches $\alpha$ at each fixed null parameter.
	Under a fixed alternative, say $b_{j,\text{alt}}$, we have
	\begin{align}
		T_n
		&= \widehat V_{jj}^{-1/2}\sqrt{n}(\widehat\beta_j-b_{j0}) \\
		&= \underbrace{\widehat V_{jj}^{-1/2}\sqrt{n}(\widehat\beta_j-b_{j,\text{alt}})}_{\dto\normaldistribution(0,1)}
		+ \underbrace{\widehat V_{jj}^{-1/2}\sqrt{n}(b_{j,\text{alt}}-b_{j0})}_{\text{diverges in absolute value}}.
	\end{align}
	Therefore, the rejection probability approaches one at each fixed alternative.
\end{remark}

\begin{remark}[Why not Standardize First?]
	One might be tempted to first standardize the full vector $\widehat\beta$ and then select the $j$-th coordinate:
	\begin{align}
		\left[\widehat V^{-1/2}\sqrt n(\widehat\beta-\beta_0)\right]_j \dto\normaldistribution(0,1).
	\end{align}
	However, notice that in this case, we are imposing the null restriction $\beta=\beta_0$ on the full vector, not just on the $j$-th coordinate.
	Since multiplying by $\widehat V^{-1/2}$ first generally mixes coefficients,
	the resulting statistic is not the same as the one obtained by first selecting the $j$-th coordinate and then standardizing.
\end{remark}

\begin{remark}[Two Cautions]\leavevmode
	\begin{enumerate}
		\item
			\emph{The normal approximation must be useful at the actual sample size.}
			A limiting result does not guarantee an accurate finite-sample approximation.
			Heavy tails or a few influential observations can make the approximation poor.
		\item
			\emph{The variance formula and its estimate must be appropriate.}
			Under heteroskedasticity, homoskedastic standard errors generally give the wrong normalization.
			With dependent observations, even the \gls{iid} heteroskedasticity-robust formula may be inappropriate.
			A poor variance estimate can also distort rejection probabilities in finite samples.
	\end{enumerate}
\end{remark}

\begin{remark}[The Exact Normal Regression Test]
	If instead we impose $e\given X\sim\normaldistribution(0,\sigma^2I_n)$,
	with $X$ of full column rank and $n>k$, then under $H_0$,
	\begin{align}
		\left.\frac{\widehat\beta_j-b_{j0}}{\sqrt{s^2[(X\T X)\inv]_{jj}}}\given[\right] X\sim t_{n-k}
		\quad\text{where}\quad
		s^2=\frac{\widehat e\T\widehat e}{n-k}.
	\end{align}
	The standardized numerator is normal and independent of
	$(n-k)s^2/\sigma^2\sim\chi_{n-k}^2$, conditional on $X$.
	Thus $t_{n-k}$ quantiles give an exact conditional test.
	And as $n\to\infty$, the $t$ distribution approaches the standard normal distribution.
	Without these normal homoskedastic assumptions,
	the so-called ``$t$-statistic'' need not have an exact $t$ distribution.
\end{remark}

\begin{remark}[Terminology: $t$-Test v.s.\ $z$-Test]
	The distinction is the reference distribution: a $t$ distribution for a $t$-test and a standard normal distribution for a $z$-test.
	A $z$-test can be exact, for example when testing a normal mean with known variance.
	Here we use a heteroskedasticity-robust statistic with an asymptotic standard normal reference distribution.
	Robustness comes from the variance estimator, not from the choice of reference distribution.
	Such regression tests are nevertheless commonly called $t$-tests.
\end{remark}

\begin{remark}[A Note on the History of the $t$-Test]
	Student's $t$-test was developed and published by William Sealy Gosset in 1908 in \emph{Biometrika} while working at the Guinness brewery.
	He published under the pseudonym ``Student'' to avoid conflicts with his employer's policy on publishing.
	The term \emph{student's $t$ distribution} was popularized by Sir Ronald A.\ Fisher in the early 20th century.
\end{remark}

\vfill
\section*{Acronyms}
\printglossaries

\end{document}
